\clearpage\begingroup\SourceLeaf{645}
\ArithmeticHead{10}{L}{H. GEPPERT, ÜBER GAUSS’ ARBEITEN ZUR MECHANIK UND POTENTIALTHEORIE.}
\renewcommand{\thefootnote}{\arabic{footnote})}\setcounter{footnote}{8}
\begin{center}\textls[100]{3. Weitere Versuche.}\end{center}
\noindent\hspace*{\parindent}War damit die Theorie der Fallversuche abgeschlossen, so blieb es noch
Sache des Experimentes, sie zu bestätigen, bezw. die bisher gefundene südliche
Abweichung durch exaktere Versuche als nicht vorhanden nachzuweisen.
Zu diesem Zwecke unternahm \textsc{Benzenberg} im folgenden Jahre 1804 eine Reihe
von Versuchen in einem verlassenen Kohlenschachte, der Alten Rosskunst bei
Schlebusch in der Mark; der Schacht wurde gewählt, um die störenden atmosphärischen
Einflüsse auszuschalten, dafür wurden die Versuche aber durch
Feuchtigkeit sehr beeinträchtigt. Die Fallhöhe betrug $85{,}1$ m, und aus 29
sehr sorgfältig angeordneten Versuchen ergab sich eine östliche Abweichung
von $11{,}5\pm2{,}9$ mm, während die Theorie dafür den Wert von $10{,}4$ mm erfordert,
und eine nördliche Abweichung von $1{,}6\pm3{,}8$ mm. Schon die letzte
Angabe zeigt, dass den Versuchen trotz aller Vorsicht noch eine grosse
Ungenauigkeit anhaftete, und in der Tat fielen die Kugeln zwischen 43 mm
nach Norden und 34 mm nach Süden, zwischen 45 mm nach Osten und
$22{,}5$ mm nach Westen, und der Gesamteindruck der Versuche ist schlechter
als der der Hamburger Experimente. Jedenfalls war aber die vorher gefundene
südliche Abweichung dadurch noch mehr in Frage gestellt\footnote{Auch diese Versuche veröffentlichte \textsc{Benzenberg} in dem oben zitierten Buche.}.

Mit mehr Geschick und Ausdauer ging 25 Jahre später \textsc{Reich} an diese
Frage heran\footnote{F. \textsc{Reich}, \textit{Fallversuche über die Umdrehung der Erde}, Freiberg 1832. Vergl. \textsc{Poggendorffs} Ann.
d. Phys. 29, 1833, S. 494—501.}. Er stellte seine Versuche im Dreibrüderschachte zu Freiberg
bei einer Fallhöhe von $158{,}54$ m an, und aus 106 sehr präzisen Versuchen
ergab sich eine östliche Abweichung von $28{,}6$ mm (theoretischer Wert $27{,}5$ mm)
und wiederum eine südliche Abweichung von $2{,}87$ mm. \textsc{Reich} selbst hat seine
Versuche falsch berechnet, die richtigen Werte gibt erst \textsc{Gauss} in einem
Briefe an \textsc{Gerling} vom 30. 12. 1852 an\footnote{Vergl. Werke XI, 1, S. 39—42. Briefwechsel zwischen \textsc{C. F. Gauss} und \textsc{Ch. L. Gerling}, herausgegeben
von \textsc{Cl. Schaefer}, Berlin 1927, S. 790—793.}; \textsc{Gauss} errechnet daselbst auch
den wahrscheinlichen Fehler dieser Abweichungen zu $3{,}5$ bezw. $3{,}4$ mm, sodass
auch hier die südliche Abweichung durchaus nicht als erwiesen angesehen
werden kann. Die Abweichungen zwischen den einzelnen Versuchs\nobreakdash-
\clearpage\SourceLeaf{646}\ArithmeticHead{11}{R}{I. NACHWEIS DER ERDROTATION.}
\HistoricalFoot{}{2\textsuperscript{*}}
\noindent ergebnissen sind sehr stark, die Kugeln fielen mit beträchtlichen Lotabweichungen
nach allen Himmelsrichtungen. Man kann nach allen diesen Versuchen \textsc{Gauss}
nur Recht geben, wenn er sagt, »dass die Fallversuche eigentlich wenig geeignet
sind die Drehungsbewegungen der Erde erkennbar zu machen, da sie
nach den kostspieligsten Zurüstungen doch immer nur höchst rohe Resultate
geben können«\footnote{Werke XI, 1, S. 39.}. In neuerer Zeit griff \textsc{Hall}\footnote{E. H. \textsc{Hall}, \textit{Do falling bodies move south?} Physical Review, 17, 1903, S. 179—190, 245—254.
\textsc{Ph. Furtwängler}, Enc. d. math. Wiss. IV, 7, S. 50 ff.} im Jahre 1902 diese Versuche
wieder auf; aus 948 Versuchen bei einer Fallhöhe von 23 m ergab sich eine
östliche Abweichung von $1{,}5\pm0{,}2$ mm (theoretisch $1{,}8$ mm) und eine südliche
Abweichung von $0{,}05\pm0{,}04$ mm; man sieht, dass auch mit den besten Mitteln
diese Versuche nicht restlos befriedigend ausfallen. Jedenfalls kann man das
Ergebnis aller dieser Versuche dahin zusammenfassen, dass die östliche Abweichung
sichergestellt ist, während innerhalb der Fehlergrenzen keine südliche
Abweichung nachgewiesen ist. Die südliche Abweichung der fallenden
Körper ist also niemals experimentell gesichert worden, trotzdem dies auch
heutigen Tages noch häufig behauptet wird. Die Suche nach ihrer Erklärung
löste in der Mitte des 19. Jahrhunderts auf Betreiben von \textsc{Oersted}\footnote{Vergl. Correspondance de H. C. \textsc{Oersted} II, Copenhague, 1920, Brief an \textsc{Herschel} vom 14. 9.
1846, S. 398 ff. Zur Geschichte der weiteren Versuche vergl. \textsc{Ferd. Rosenberger}, \textit{Die Geschichte der
Physik}. Braunschweig 1887—1890, 3, S. 432—437; P. \textsc{Stäckel}, Enc. d. math. Wiss. IV, 6, S. 485.}
eine lebhafte Diskussion aus, die jedoch, ebenso wie die hieran anschliessenden
Versuche, wertlos ist. Hingegen verdient ein anderer Punkt bei
den vorangehenden Experimenten Beachtung. Alle Beobachter konstatierten,
dass selbst unter grösster Vorsicht beim Ablassen die Kugeln während des Falles
eine Rotation erleiden, und zwar eine verhältnismässig grosse um eine horizontale
Achse und eine geringere Drehung um die vertikale Achse. Dieses
Phänomen, das im übrigen keinen Einfluss auf die Versuche gehabt zu haben
scheint, wurde von den genannten Beobachtern nicht erklärt, und erst 1928
von \textsc{Denizot}\footnote{A. \textsc{Denizot}, \textit{Sur un phénomène observé par Guglielmini à Bologne en 1791.} Atti del Congresso
internaz. dei Matematici a Bologna 1928, 6, S. 475—482.} theoretisch zu begründen versucht; die Rotation der fallenden
Kugeln findet nur auf der rotierenden Erde statt und scheint also in gleicher
Weise ein Beweis für deren Achsendrehung zu sein, wie es die östliche Abweichung
bei den angestellten Versuchen sein sollte.
\clearpage\endgroup
